\section{\heiti 讨论}
\label{sec:discussion}

\subsection{\heiti 循环位置与容量增长}
\label{sec:loop-side}

第~\ref{sec:compare}~节的对照可以归结为一个位置选择。ALIT 的循环在编码侧：每轮都把图像与既有的隐 token 一并重新编码，并按上一轮的重建质量掩掉已蒸馏充分的部分；本文的循环在解码侧：编码器只前向一次，每步只运行读出块 $\Phi$。这一取舍把编码器的开销与读出步数解耦，使任意预训练编码器都可以直接接入，而不必为每个 token 预算重新编码；代价是编码器不具备“因局部区域重建不足而重新组织输入”的通路。

位置选择还决定了容量如何增长。ALIT 每轮拼接新的可学习 token，新 token 构成新的记忆槽，可在后续轮次中专门承载某一区域，这是它报告 token 专业化与目标发现的来源。本文的容量增长是确定性的分块窗口：token 集合固定为 $K$，第 $t$ 步仅增加可读入的 token 数。由此得到两点性质：不引入额外参数；第 $t$ 步的码与完整前向中第 $t$ 步逐位相同（式~\eqref{eq:recur}），因此停在任意一步的结果是精确的。代价是第~\ref{sec:complexity}~节观察到的复杂度效应无法转化为逐图的预算差异——所有图像共用同一条逐步曲线，仅高度不同。

\subsection{\heiti 预算的取用与匹配}
\label{sec:deploy}

第~\ref{sec:complexity}~节表明，复杂图像在每一项容量档上的误差都更高，容量增加缩小的是复杂度惩罚的绝对幅度：最高与最低三分位的 L1 差从 63M 档的 5.655 / 4.992 / 2.811 px 降到 240M 档的 3.629 / 2.821 / 1.504 px，但相对差基本不变（240M 档仍为 66.7\% / 49.7\% / 22.5\%）。这说明“不同图像对预算的需求确实不同”的动机成立，但相对结构不随容量改变；同一条逐步曲线的形状对所有图像相同，逐图差异只体现在高度上。

因此本文的建议是：\textbf{由部署方按场景精度预设预算}，而不是让模型在推理时逐图预测预算。例如低精度场景取 2--8 个 token，高精度场景取 16--32 个 token。这一建议与“预算可自适应调整”的设计一致：用户只需训练一个模型，就能在多个预设预算下取用。若要进一步按图分配，需要额外的判断通路，本文未做。

第~\ref{sec:ksweep}~节还给出两条可操作的边界：一是容量，在本协议下 240M 档（合计 481.5 M）是甜点，继续加宽到 721M 反而退化；二是步数与预算的匹配，主实验取 $K=2T$（每步 2 个 token），使每一步都有实质的信息增量。

\subsection{\heiti 同类机制的异同}
\label{sec:alit}

第~\ref{sec:compare}~节只把逐块 PCA 并入按预算分层的同值数对照（其他分词器的口径差异见该节说明），本文与 ALIT 的机制差别集中在三处。

\textbf{循环放在哪一侧。} ALIT 的循环在编码侧，本文在解码侧，取舍见第~\ref{sec:loop-side}~节。

\textbf{容量如何增长。} ALIT 每轮新增记忆槽，本文用读窗口扩张；前者可让 token 分工，后者不增参数，且停在任意一步的结果精确。

\textbf{码是否量化。} ALIT 量化到共享的一维码本，“长度”同时是比特预算，可与 VQGAN、TiTok 在 FID 与 bpp 上比较；本文不做量化、不建熵模型，压缩率只以 token 数度量，因此第~\ref{sec:compare}~节的 PSNR 对照只在本协议内成立，不能换算成比特率；与其他分词器的数值比较需在统一协议下复现其训练与评测后再做。

这三点共同说明：两者都在“权重共享的循环 + 每步容量递增”这一框架下工作，但一个把容量放在写侧，一个放在读侧，因而得到的性质不同——前者可以按图分配长度，后者把预算的决定放在推理期。

\subsection{\heiti 适用边界}
\label{sec:limitation}

本文的结果有若干边界。第一，受控模型为 481.54 M 参数（编码器与解码器各约 240 M）、在 ImageNet-100 上从零训练，其绝对重建质量不能外推到大规模预训练模型；本文只在同一协议内比较，不对绝对数值作跨协议推广。第二，本文不做量化与熵编码，“压缩率”以 token 数度量，不是真实比特率；本文的模型底座（无 VQGAN、无对抗损失）、训练数据与分辨率也与 TiTok、ALIT 不同，因此表~\ref{tab:budget} 不与其并列数值，待按统一协议复现后再比。第三，评价指标以像素重建为主，没有下游任务；压缩表示对图像描述、视觉问答等任务是否够用仍需另行检验。第四，单步与循环的训练预算未对齐（阶段 1 为 45 epoch 从零，阶段 2 为 6 epoch 热启动），本文不主张“循环优于单步”；同预算对照（从同一单步权重分别跑“继续单步”与“BPTT”）是下一步工作。第五，复杂度分析使用分布外的 SAVOIAS 数据，结论只说明“在该人评复杂度轴上呈对齐趋势”，不能推断 ImageNet-100 验证集内部的复杂度分布。第六，第~\ref{sec:ksweep}~节的容量阶梯在 360M 档出现反转，原因（同一学习率与同一日程下宽模型更难优化）尚未与数据侧因素分离；旧协议（63M/K=64）下 $K>64$ 的退化同理未完全解耦。第七，本文只检验了 ALIT 所列四条轴中的一条（复杂度），没有检验“对训练分布的熟悉度”，也没有分析“token 专业化”；对这两点本文既不主张也不否定。第八，按预算分层的对照（表~\ref{tab:budget}）表明：预算越紧本文的码越占优，预算越松逐块线性码越占优；本文的主预算点（39{,}936 值）并非同值数最优（27.94 dB 对 PCA 的 35.75 dB），因此本文不把“同值数下最优”作为主张。

\subsection{\heiti 时序的延伸}
\label{sec:video}

读出本来就是一条时间轴：第 $t$ 步在第 $t-1$ 步的状态上再加一步。单帧里它对应“同样的内容看得更全”，放到视频上则对应“时间由过去到现在”。这是本文最看好的一条延伸：在视频里，循环不只是把一步拆成几步，而是让前一帧的结果成为后一帧的起点。

视频的相邻两帧几乎总是相似的：机位不动时背景逐像素不变，只有物体和光照在动；手持拍摄时差异也只集中在画面的一小部分。后一帧的大部分内容前一帧已经解释过，从零重建等于把同一块背景的 token 预算再花一遍。

一种自然的用法是：每帧只走少量步数，第一步的初值取上一帧的读出状态，由它承担“没变的部分”，新增的几步只补“变了的部分”。帧间差异小就少走几步，运动大就多走几步，镜头切换时退回全量读出。第~\ref{sec:loop-side}~节的两点性质在这里正好合用：提前停止是精确的，少走几步不会改动已读出的结果；循环在解码侧，跨帧搬过去的只是一个读出状态，不必把上一帧重新编码一遍。

这种复用与视频编码器的帧间预测不在同一层：后者在比特层压掉静止背景，管的是码率；前者在特征层复用状态，省的是 token 预算。

真要走到这一步，还差三件事。一是训练：监督只在单帧内由粗到细，帧间状态没进过训练目标，需要把相邻帧接成序列、跨帧传状态后再微调。二是误差累积：上一帧的偏差会带进下一帧，系统要周期性插关键帧刷新。三是评价：视频还要看时序一致性，本文没有这项指标与视频数据。
